\begin{center}
\zihao{4}\sffamily 摘\quad 要
\end{center}

大模型智能体不只是回答问题，还会读取检索内容和持久化记忆，规划任务，并调用邮件、文件、网络和代码工具。攻击者可以把恶意指令藏在这些外部内容中，诱导智能体把一段文字变成真实动作，造成内部邮件外送、受保护文件读取、越权 API 调用或危险数据写入。只检查用户提示词或模型生成的文字，不能单独保证即将执行的工具参数安全；一旦动作已经发生，信息外送和文件写入往往无法撤回。

本作品设计并实现\textbf{灵鉴Agentscope 沙箱安全引擎}，一套可复用于不同智能体框架的\textbf{运行时安全中间件}。它不替换原有模型和工具，而是在真实动作发生前提供独立安全屏障：确定性规则先拦截明确高危，未解决信号再进入本地模型，仍需升级时先在本地完成确定性脱敏，再把最小必要信息交给外部裁判；所有路径统一归约为\textbf{允许、询问、拒绝、脱敏后允许}四种动作。作品已接入 OpenClaw，在运行前、模型输出交付前、工具执行前和消息交付前四个真实执行点形成连续监督链。

作品构造并封印了\textbf{300 条结构化评估夹具}（整体英文 150 条、中文 150 条；风险 180 条、安全 120 条），覆盖九类风险；其中 54 条第一方对抗变换均为英文。作品最终交付的是一套能够在动作发生前判定、阻断、留证并复核的智能体运行时安全中间件。

\noindent\textbf{关键词}：大模型智能体；运行时安全中间件；工具调用监督；三级级联；确定性脱敏；可复现评测
